\section{Background and Motivation}

\subsection{Streaming Graph Processing}

A graph $G=(V,E)$ consists of a vertex set $V$ and an edge set $E$ representing relationships between vertices. Many graph algorithms, including breadth-first search (BFS), single-source shortest paths (SSSP), and connected components (CC), iteratively update vertex values using a vertex-centric processing model. In each iteration, only a finite set of active vertices in the current worklist is processed, and vertices that require further processing are added to the next worklist until convergence [Ligra].

A streaming graph evolves through a continuous sequence of updates, typically grouped into batches of edge insertions and deletions. Each batch typically modifies only a small portion of the graph, and applying batch $\Delta G_t$ transforms $G_t$ into $G_{t+1}$. Systems such as STINGER, GraphOne, and Terrace maintain mutable graph representations to support topology updates and graph analytics [STINGER, GraphOne, Terrace]. To maintain up-to-date query results, a straightforward approach reruns the graph algorithm after each batch update, obtaining $R_{t+1}=\mathcal{A}(G_{t+1})$ even when only a small portion of the topology has changed.

\subsection{GPU Graph Processing beyond Device Memory}

The massive parallelism of GPUs makes them a natural platform for vertex-centric graph processing. Early frameworks map this computation to GPU threads and organize parallel traversal and value propagation [Medusa, CuSha, Gunrock]. However, their scalability is constrained by device memory capacity, as large graphs can exceed the memory available on a single GPU and require out-of-memory (OOM) processing.
% TODO: Add the original Medusa, CuSha, and Gunrock papers/bibliography entries.

For static graphs, OOM systems retain graph data in host memory and transfer the portions needed for GPU computation. Subway constructs and transfers active subgraphs, EMOGI enables fine-grained zero-copy reads of host-resident edges, and HyTGraph selects between explicit transfers and zero-copy access [Subway, EMOGI, HyTGraph]. CGgraph further combines a reusable GPU-resident subgraph with CPU-GPU task allocation for cooperative execution [CGgraph]. These systems avoid loading the entire graph onto the GPU; however, since PCIe transfer speeds are significantly lower than memory access speeds, reading data from the host still incurs overhead.

The dynamic graph processing system introduces the requirement of maintaining adjacency storage as updates arrive. GPU structures such as Hornet and faimGraph support topology changes within device memory [Hornet, faimGraph]. Extending GPU processing to OOM dynamic graphs requires both access to host-resident topology and efficient maintenance of graph data and index views on the GPU after updates. In contrast to static graph storage, the host representation and any GPU-resident adjacency data must remain consistent with the current update batch before computation proceeds. In streaming graphs, even though only a portion of the graph data is updated, recalculating the results after every update would significantly increase communication overhead. Therefore, merely addressing capacity constraints does not eliminate the overhead associated with continuous graph analysis.

\subsection{Incremental Processing and CPU-GPU Cooperation}

Incremental graph processing reuses previous results to avoid recomputing vertices unaffected by streaming graph update [KickStarter, GraphBolt, DZiG]. For monotonic path-based algorithms such as BFS and SSSP, a key approach records each vertex's \emph{dependency}: the predecessor that supplied its current value. When an edge deletion breaks a dependency, the algorithm traces dependent vertices, invalidates their values, and repairs them using remaining incoming edges. Inserted edges can introduce better values, which propagate through an active work list until convergence. The vertices requiring recomputation are determined by the propagation process triggered by the update, rather than by the updated edges alone. Repair may also require values from unaffected predecessors that provide alternative paths.

Grapin brings dependency-based incremental computation to GPU streaming graph processing [Grapin]. It maintains dynamic topology in host memory, performs incremental analysis on the GPU, and uses zero-copy access to support graphs exceeding device memory. It also caches frequently accessed subgraphs on the GPU to reduce repeated transfers. Nevertheless, the locality of streaming updates is not fully exploited: incremental computation limits result repair to affected vertices, but graph traversal and data structure maintenance can still scan the entire graph or process substantially more vertices and edges than required by the update. Such processing adds redundant computation and, when unnecessary adjacency data are fetched from host memory, PCIe traffic. In this execution model, the CPU primarily receives and applies graph updates and schedules GPU execution. Its otherwise idle computing resources can be used to prepare the host-resident graph data needed by affected vertices before GPU processing.

CPU-GPU cooperation offers a way to reduce such communication. CGCGraph processes shared graph data on the GPU and assigns snapshot-specific work to the CPU, reducing transfers across concurrent snapshot jobs [CGCGraph]. Related work, POEGA, uses compact proxy graphs and concurrent GPU analysis to reduce OOM accesses across a sequence of snapshots [POEGA]. Both target analysis of multiple available snapshots; their gains rely on sharing or amortizing work across those snapshots. These execution models do not address online streaming computation, which must maintain the latest result as each new update batch arrives.

These observations motivate our CPU-GPU cooperative incremental processing of OOM streaming graphs. By using otherwise idle CPU resources for data preparation and further exploiting update locality in graph traversal and data structure maintenance, we aim to reduce redundant computation and PCIe traffic.
